\chapter{نمونه اجرای گام‌به‌گام روش‌ها}\label{app:step-by-step}

این پیوست، اجرای یک نمونه مشترک از پنج روش را ثبت می‌کند. شناسه نمونه، برچسب واقعی، متن ورودی، تصمیم‌های میانی، زمینه‌های انتخاب‌شده و پاسخ نهایی گزارش شده‌اند.

\section{نمونه مشترک}
\begin{latin}\footnotesize\begin{verbatim}
sample_id: dr-test-029c26d1b1461025
gold label: yes
post excerpt: “I get more high strung, anxious and depressed every day ...
                Living is becoming a chore.”
\end{verbatim}
\end{latin}

\section{روش \lr{Vanilla RAG}}
\begin{latin}\footnotesize\begin{verbatim}
pipeline:
passage_retrieval -> reranking -> answer_generation
passage_retrieval: 50 candidates
reranking: top 10 candidates
answer input: top 5 evidence passages
P1: “Seemed sad or depressed for several hours?”
P2: “Feeling down, depressed, or hopeless?”
P3/P5: clinical depression examples
answer: yes
cited_passage_ids: [P1, P2]
\end{verbatim}
\end{latin}
توضیح روش بر اساس دو قطعه نخست است: متن نمونه به‌صراحت از افسردگی روزانه و دشوارشدن زندگی سخن می‌گوید و قطعه‌های بازیابی‌شده این توصیف را با خلق افسرده و ناامیدی مرتبط می‌کنند.

\section{روش \lr{HippoRAG 2}}
\begin{latin}\footnotesize\begin{verbatim}
pipeline:
fact_generation -> fact_filtering -> fact_retrieval -> seed_weighting
-> PPR -> passage_ranking -> answer_generation
queries:
anxiety --can co-occur with--> depression
feeling high strung --is associated with--> anxiety
depressive mood --is a symptom of--> depression
living is a chore --may indicate--> depression
retrieved fact candidates: 20
selected facts: 12
initial passages: 6
ranked passage candidates: 20
answer: yes
cited_passage_ids: []
evidence_passage_ids: [P1, P2, P3, P4, P5]
\end{verbatim}
\end{latin}
در این نمونه، واقعیت‌های استخراج‌شده به‌صورت سه‌تایی تولید شدند. پس از پالایش و بازیابی، قطعه‌های مرتبط با خلق افسرده و ناامیدی در رتبه‌های بالاتر قرار گرفتند و پاسخ نهایی «بله» بود.

\section{روش زنجیره فکر \lr{CoT}}
\begin{latin}\footnotesize\begin{verbatim}
pipeline:
question + valid_labels -> cot_generation -> answer
retrieval: none
graph traversal: none
LLM calls: 1
answer: yes
\end{verbatim}
\end{latin}
مدل بر پایه مشاهده مستقیم «افسردگی هر روزه» و دشوارشدن زندگی، برچسب «بله» را انتخاب کرد؛ در این روش هیچ بازیابی متن یا پیمایش گرافی انجام نشد.

\section{روش \lr{ToG-1}}
\begin{latin}\footnotesize\begin{verbatim}
pipeline:
topic_entity_extraction -> graph traversal -> answer_generation
topic entity: depressed
retrieved triplet: (Depressed, manifests_with, negative mood)
answer prompt: original post + triplet + allowed labels
answer: yes
\end{verbatim}
\end{latin}
موجودیت «افسرده» از متن استخراج شد و سه‌تایی مرتبط با خلق منفی از گراف بازیابی گردید. مدل با ترکیب متن اصلی و مسیر گرافی، برچسب «بله» را تولید کرد.

\section{روش \lr{ToG-2}}
\begin{latin}\footnotesize\begin{verbatim}
pipeline:
TopicEntityExtraction -> TopicPrune -> CombinedRelationPrune
-> graph traversal -> context ranking -> SufficiencyAndAnswer
entities: [high strung, anxious, depressed, depression]
seed entities: 10
PPR iterations: 26
calls: 1_extract, 2_topic_prune, 3_relation_prune, 4_sufficiency
ranked contexts:
  chunk 6596f07f: marked depressed mood
  chunk 3e47f1d7: hopeless daily / severe depression assessment
answer: yes
stopping_reason: llm_sufficient
selected paths: none in this sample
\end{verbatim}
\end{latin}
در این نمونه، پس از استخراج و هرس موجودیت‌ها، قطعه‌های مرتبط رتبه‌بندی شدند. مدل زبانی کفایت زمینه را تأیید کرد و حلقه با پاسخ «بله» پایان یافت.

\section{نمونه استخراج ادعا از توضیح \lr{ToG-2}}
\begin{latin}\scriptsize\begin{verbatim}
case claims:
MCASE_1: the poster -- reports -> depressed every day
MCASE_2: the poster -- reports -> living is becoming a chore
clinical claims:
MCLINICAL_1: depressed every day -- aligns with -> marked depressed mood
MCLINICAL_2: living is a chore -- aligns with -> marked depressed mood
MCLINICAL_3: depressed every day -- aligns with -> hopeless daily
MCLINICAL_4: living is a chore -- aligns with -> hopeless daily
\end{verbatim}
\end{latin}
در ارزیابی توضیح، ادعاهای مربوط به شخص از ادعاهای بالینی جدا شدند و فقط ادعاهای بالینی با فکت‌های مجموعه دانش مقایسه شدند.
